\subsection{A zero-cost collar can have a signed interim cost reserve}
\label{aligned-zero-cost-collar}

Equal option premiums can conceal a changing net time value. A purchaser may
pay for a call that limits price increases and finance it by writing a put that
gives up price decreases below a lower strike. Even when those premiums cancel,
the two options can have different time values before expiry. IFRS~9.B6.5.31
expressly includes a purchased/written option combination with net nil time
value at designation: its related time-value changes enter other comprehensive
income, although the cumulative change over the full relationship is zero
\citep[IFRS~9.B6.5.31--32]{eu20231803}.

Consider a highly probable purchase of ten units. A purchased call has strike
11 and a written put has strike 9, both expiring on the purchase date. Spot is
initially 10, so both intrinsic values are zero. The call premium paid and put
premium received are each 6. They are two gross cash flows and two actual
contract balances, even though their net cash effect is zero. Assume external
contracts, a qualifying combined designation that is not a net written option,
documented matching risk components, complete values and continuing
qualification. IFRS~9.6.2.5--6 permits joint designation subject to those
eligibility restrictions; equal premiums alone do not establish them
\citep[IFRS~9.6.2.1--6, 6.3.7, 6.4.1]{eu20231803}.

Let \(q\) be the common quantity, \(L<U\) the lower put and upper call
strikes, and \(S\) the spot price at a measurement date. Denote the positive
call asset by \(C\) and positive put liability by \(P\). Their supplied
nonnegative time values are \(T_C\) and \(T_P\). The contract payoffs give
\[
 \begin{aligned}
 I_C&=q\max(S-U,0),& C&=I_C+T_C,\\
 I_P&=q\max(L-S,0),& P&=I_P+T_P.
 \end{aligned}
\]
The put enters the purchaser's net position with a negative sign. In this
finite measurement profile, zero discount and the supplied matching item-risk
measure make the hedged-item change exactly opposite to \(I_C-I_P\).
No carry, credit or basis difference remains. The lower-of test therefore gives
the signed intrinsic reserve \(R_I=I_C-I_P\), with no ineffectiveness, while
the fully aligned cost reserve is \(R_T=T_C-T_P\). Consequently
\[
 C-P=R_I+R_T.
\]
The reserves retain different origins. A zero sum cannot justify dropping
either one. This use of the matching item measure follows the declared
designation; a hypothetical derivative cannot import instrument features
absent from the item \citep[IFRS~9.6.5.11, B6.5.5]{eu20231803}.

One intermediate valuation has spot 11.7, call asset 15 and put liability 3.
Call intrinsic value is 7, leaving call time value 8; the put has no intrinsic
value and time value 3. Thus \(R_I=7\) and \(R_T=5\). At a later valuation,
spot is 8.3, the call asset is 3 and the put liability is 15. Put intrinsic
value is now 7 and its time value is 8, giving \(R_I=-7\) and
\(R_T=-5\). The earlier common-loss restriction for purchased-call time
values cannot describe this signed reversal.

These synthetic marks also have a simple consistency check. At zero discount,
terminal-price weights on \((6,9,10,11,14)\) of
\((1/10,1/10,1/10,1/5,1/2)\) give mean price 11.7, mean call payoff 15
and mean put payoff 3 for ten units. Reversing the weights gives mean price
8.3 and option values 3 and 15. Equal weights give the inception price 10 and
premiums 6 each. These are separate mathematical witnesses for the supplied
marks; they are neither estimated probabilities nor a joint transition or
forecast model across dates.

At expiry, both time values vanish. The purchaser pays the supplier \(qS\),
receives call payoff \(I_C\) and pays put payoff \(I_P\), as distinct cash
flows. The final inventory cost, after the intrinsic reserve is included in
its basis, is
\[
 B=qS-I_C+I_P=q\min\!\bigl(\max(S,L),U\bigr).
\]
Below 9 the written put raises the ten-unit cost to 90; between 9 and 11 the
cost is the supplier price; above 11 the call limits it to 110. Terminal prices
6, 9, 10, 11 and 14 therefore give costs 90, 90, 100, 110 and 110. With initial
capital 300, cash closes at 210, 210, 200, 190 and 190. Derivative balances,
both reserves and earnings close at zero, so cash plus inventory remains 300.
The time-value basis adjustment is zero because \(R_T=0\) at expiry, exactly
as B6.5.31(a) requires. Its interim OCI history remains part of the account.

\texttt{AlignedZeroCostCollar/1} returns the two actual contracts, both signed
reserves, proposed balanced journals and all five gross cash origins. It
requires sufficient cash at each supplied payment step: call premium before
put premium receipt, then supplier payment before derivative settlements.
A future observation, missing qualification, wrong component value or netted
supplier amount returns no partial report. The method supplies no option
pricing model, tax treatment or cash-flow-activity classification. Unmatched
actual/aligned changes under B6.5.33 and reserve attribution to heterogeneous
stopped purchase portions require separate support.
